\part{Práctica 2: Modelado y Gestión}

\section{Recursos y Organización del Equipo}
Para este proyecto se utilizará una arquitectura basada en tecnologías web modernas, apoyada por asistentes de IA (OpenCode, Engram) para acelerar el desarrollo. La gestión del ciclo de vida del software se realizará en GitHub mediante la metodología ágil SCRUMBAN \cite{SCRUMBAN}.

El repositorio oficial del proyecto puede consultarse en el siguiente enlace: \href{https://github.com/Jsoriano99/qr-restaurante-iw}{Repositorio GitHub - qr-restaurante-iw}

En cuanto a la organización del equipo, al tratarse de un desarrollo individual, el único integrante asumirá los siguientes roles simultáneamente para asegurar el cumplimiento del ciclo de vida del software:

\begin{description}
    \item[Scrum Master / Pruebas:] Gestión del tablero Kanban en GitHub y validación de calidad.
    \item[Modelado y Diseño:] Especificación UWE (MagicDraw / Draw.io).
    \item[Desarrollo Front-end \& Back-end:] Implementación del sistema web y APIs de códigos QR.
    \item[Gestor de la Base de Datos:] Diseño y mantenimiento del historial de afluencias y usuarios.
\end{description}

\section{Modelado de Requisitos en UWE}
La metodología UWE (UML-based Web Engineering) distingue fundamentalmente entre dos tipos de casos de uso \cite{UWE2014}:
\begin{itemize}
    \item \textbf{Acciones de Navegación (\textit{«browsing»}):} Casos de uso donde el usuario simplemente interactúa con la información, consultando datos sin alterar el estado del sistema. 
    \item \textbf{Acciones de Proceso (\textit{«processing»}):} Casos de uso que implican transacciones o modificaciones sustanciales en la lógica de negocio o en la base de datos.
\end{itemize}

\subsection{Diagrama de Casos de Uso}
En nuestro modelo, podemos observar claramente esta distinción: el caso de uso de \textbf{Listado de Restaurantes} se clasifica como \textit{«browsing»}, ya que es una consulta de disponibilidad. Por otro lado, el \textbf{Registro de Entrada} es eminentemente \textit{«processing»}, puesto que altera el aforo en tiempo real.

\begin{figure}[hbt!]
    \centering
    \includegraphics[width=0.8\textwidth]{images/caso_de_uso.jpeg}
    \caption{Diagrama de Casos de Uso - Operaciones de navegación y proceso.}
    \label{fig:casosuso}
\end{figure}

\subsection{Diagramas de Actividades}
Para comprender mejor la lógica de negocio y las transacciones dentro del sistema, se han modelado los diagramas de actividad de los procesos clave de la plataforma. En ellos se detallan los flujos de ejecución, validaciones y puntos de decisión.

A continuación, se presentan los diagramas de actividad correspondientes al registro del restaurante, la generación del QR, el escaneo de entrada y el registro de salida.

Cabe destacar el flujo correspondiente al registro de entrada (Figura \ref{fig:escaneo_entrada}), donde se modela la lógica de validación de aforo, incluyendo las acciones del usuario (\textit{«userAction»}) y las acciones del sistema (\textit{«systemAction»}). Adicionalmente, se hace uso del \textit{interactionPin} para inyectar los datos relativos al número de acompañantes en el flujo principal.

\begin{figure}[hbt!]
    \centering
    \includegraphics[width=0.7\textwidth,height=0.4\textheight,keepaspectratio]{images/Restaurante-Registration.png}
    \caption{Diagrama de Actividades - Registro de Restaurante.}
    \label{fig:registro_restaurante}
\end{figure}

\begin{figure}[hbt!]
    \centering
    \includegraphics[width=0.7\textwidth,height=0.4\textheight,keepaspectratio]{images/generacion-qr.png}
    \caption{Diagrama de Actividades - Generación de código QR.}
    \label{fig:generacion_qr}
\end{figure}

\begin{figure}[hbt!]
    \centering
    \includegraphics[width=0.7\textwidth,height=0.4\textheight,keepaspectratio]{images/escaneo-entrada.png}
    \caption{Diagrama de Actividades - Escaneo y Registro de Entrada (HU-03).}
    \label{fig:escaneo_entrada}
\end{figure}

\begin{figure}[hbt!]
    \centering
    \includegraphics[width=0.7\textwidth,height=0.4\textheight,keepaspectratio]{images/registro-salida.png}
    \caption{Diagrama de Actividades - Registro de Salida.}
    \label{fig:registro_salida}
\end{figure}

\clearpage

\section{Metodología de Gestión}
El proyecto se ha gestionado empleando la metodología ágil multihíbrida \textbf{SCRUMBAN} \cite{SCRUMBAN}, utilizando para ello la herramienta \textit{GitHub Projects}. Esta aproximación permite una visualización eficiente del flujo de trabajo continuo (Kanban) con iteraciones estructuradas (Scrum).

Las Historias de Usuario (HU) del proyecto han sido modeladas como \textit{Issues} individuales dentro del repositorio. Cada issue cuenta con metadatos específicos, tales como etiquetas de \textit{story points} (por ejemplo, \texttt{points:3} o \texttt{points:5}) para medir su complejidad y esfuerzo de desarrollo, y etiquetas de prioridad (como \texttt{priority:high}) para organizar el backlog.
